\section{Introduction}\label{sec:intro}

Modern JavaScript applications are assembled from hundreds of npm packages. Each package declares its
own dependencies through SemVer ranges, and package managers resolve and install these transitively.
This creates a dense, directed web of trust. A compromise of a single widely used package can reach a
large part of the ecosystem~\citep{zimmermann2019smallworld,kikas2017structure,decan2019empirical}.
The exposure is no longer hypothetical. Account takeovers and malicious maintainer handovers have
repeatedly delivered malware through trusted packages~\citep{npm2018eventstream,github2021uaparser,
aikido2025chalkdebug,ohm2020backstabber}.

In 2025 the threat changed qualitatively, as worms began to propagate themselves:
\begin{itemize}
  \item \textbf{Shai-Hulud} (September) ran a malicious script at install time. The script harvested
        npm publish tokens from developer machines and CI runners, then used each stolen token to
        republish up to \CalOneCap{} of the victim's own packages with the same payload
        injected~\citep{stepsecurity2025shaihulud,socket2025shaihulud,cisa2025npm}.
  \item \textbf{Shai-Hulud~2.0} (November) used a preinstall hook and raised the per-token cap to
        \CalTwoCap{} packages~\citep{wiz2025shaihulud2,datadog2025shaihulud2}.
\end{itemize}
Each newly infected package infected whoever installed it next. Spread was limited only by detection and
registry intervention. These incidents raise network-scientific questions that static vulnerability
reachability studies~\citep{decan2018vulnerabilities,liu2022demystifying,zerouali2022vulnerabilities}
cannot answer:
\begin{itemize}
  \item How do the topology of the dependency graph and the ownership structure of maintainer accounts
        jointly determine how far and how fast such a worm spreads?
  \item Which ecosystem-level defences change the epidemic threshold?
  \item Under realistic detection delays, should a limited response capacity go to patching what is
        already infected, or to hardening what is not?
\end{itemize}

We address these questions with an SIR epidemic model of the npm ecosystem that has two contagion
channels:
\begin{itemize}
  \item a \emph{dependency channel}, in which a compromised package exposes its dependents with a
        probability that depends on the dependents' SemVer constraints and on whether they execute
        lifecycle scripts;
  \item a \emph{credential channel}, in which a compromised package's maintainer token is stolen and
        the worm republishes a capped number of that maintainer's packages.
\end{itemize}

We study three research questions:
\begin{description}
  \item[RQ1 (origin topology).] How much do outbreak size and speed depend on where the worm starts:
        a hub, a bridge, a mid-tier or a leaf package?
  \item[RQ2 (ecosystem regimes).] How do strict lockfiles, disabled lifecycle scripts and trusted
        publishing shift the epidemic threshold?
  \item[RQ3 (interventions).] At a fixed response budget, how does reactive patching compare with
        proactive immunization under realistic detection delays?
\end{description}

\paragraph{Contributions.}
\begin{enumerate}
  \item \textbf{Model.} We give, to our knowledge, the first epidemic-dynamics model of self-propagating
        npm worms that couples a SemVer-weighted dependency channel with a stolen-token credential
        channel (\cref{sec:model}).
  \item \textbf{Data.} We provide a cleaned \MainN-package npm dependency graph, built from the
        registry with provenance, install-script and maintainer metadata, plus a full-closure variant
        and a synthetic benchmark (\cref{sec:data}). Cleaning removes a reward-farming spam ring that
        on its own would have dominated the spectral analysis.
  \item \textbf{Calibration and validation.} We calibrate the model to the timelines of both Shai-Hulud
        waves and validate its channel structure against the real lists of compromised packages
        (\cref{sec:calibration}).
  \item \textbf{Evidence on defences.} We show that credential compromise, not dependency topology,
        governs the reach of a worm (RQ1). Dependency-side hardening raises the epidemic threshold only
        while tokens cannot be stolen (RQ2). Proactive immunization of high-degree packages
        outperforms reactive patching once detection takes between 12 and 24 hours or more (RQ3,
        \cref{sec:results}). Null models, synthetic ensembles, sample sizes and parameter sweeps test
        the robustness of these findings (\cref{sec:robustness}).
\end{enumerate}
Code, data snapshots and the frozen experiment configuration are released with the paper.
